% ---------------------------------------------------------------------------
\section{Conclusion}
% ---------------------------------------------------------------------------

We introduced ACE, an analytical, training-free DCT rate proxy derived from the HEVC specification.
Three theoretical results establish convergence, near-optimality under the Markov-1 image model, and gradient consistency with the true HEVC rate.
In turn, establishing ACE as a reliable proxy to perform compression-aware gradient-based optimization.

We demonstrated this utility on four tasks.
At inference time, minimizing ACE with an $\ell_2$ fidelity anchor recovers $32\%$--$45\%$ of the clean--adversarial accuracy gap on without any knowledge of the attack or classifier, compared to $3\%$--$20\%$ for JPEG compression at matched quality.
At training time, applying ACE gradient steps as data augmentation (ACEAug) improves robustness to CIFAR-10-C corruption with negligible clean-accuracy cost.
In particular, ACEAug provides an additive gain on top of the strong AugMix baseline, confirming that the frequency-domain signal is orthogonal to the geometric and photometric perturbations covered by AugMix.
The same signal generalizes in three further directions.
Used to constrain an attack's own perturbation rather than to remove one, as a zero-parameter penalty on intermediate features, as a diffusion guidance term.

ACE is a few-line, parameter-free component that can be dropped into any pipeline that requires a differentiable approximation of HEVC bitrate.
Currently, the proxy targets $8\times8$ blocks; it is left to future work to explore modeling multi-scale partitioning and multiple intra-prediction modes in HEVC, and the potential applications this unlocks.